\chapter{System Architecture}\label{chap:background}

\section{Repository Overview}\label{sec:repo-overview}
The system is distributed across six repositories under the \texttt{Ejada-A}
GitHub organization. Table~\ref{tab:repo-overview} summarises the responsibility
and structure of each one.

\begin{table}[H]
\centering
\caption{Repositories and responsibilities.}
\label{tab:repo-overview}
\small
\begin{tabularx}{\textwidth}{>{\raggedright\arraybackslash}p{3.2cm}
                >{\raggedright\arraybackslash}X}
\toprule
\textbf{Repository} & \textbf{Contents} \\
\midrule
\texttt{auth-\allowbreak service} & One entry point, \texttt{server.ts} (206 lines), plus \texttt{user.model.ts} and \texttt{admin.model.ts}. There are no separate routes/controllers/services directories -- route handlers live directly in \texttt{server.ts}. \\
\texttt{products-\allowbreak service} & Same shape: \texttt{server.ts} (191 lines), \texttt{product.model.ts}, \texttt{category.model.ts}. \\
\texttt{orders-\allowbreak service} & \texttt{server.ts} (128 lines), \texttt{order.model.ts}, and its own copy of \texttt{product.model.ts}. \\
\texttt{payments-\allowbreak service} & \texttt{server.ts} (98 lines), its own copies of \texttt{order.model.ts} and \texttt{product.model.ts}. \\
\texttt{ecomm-ui} & A Next.js 15 (App Router) application: storefront pages, an \texttt{/admin} section, and a set of \texttt{/api/*} routes that act as a backend-for-frontend (BFF), proxying to the four backend services over HTTP. \\
\texttt{Terraform-\allowbreak code} & The Terraform \texttt{network} and \texttt{oke} modules, the Helm chart (\texttt{helm/ecommerce-app}), a parallel set of raw Kubernetes manifests (\texttt{kubernetes/}), ArgoCD Application/Ingress definitions (\texttt{argocd/}), and the deployment runbook \texttt{DEPLOYMENT.md}. \\
\bottomrule
\end{tabularx}
\end{table}

Each backend service uses a flat Express structure. Routes and request handlers are
defined in \texttt{server.ts}, where each handler calls the corresponding Mongoose
model directly; there is no separate controller, service, or repository layer.

\begin{infobox}[Note]
Automated tests are not currently included in the six repositories.
\texttt{Terraform-code}'s \texttt{test-runner.yml} workflow contains a placeholder
check (\texttt{echo "Runner works"}) rather than an application or infrastructure
test suite.
\end{infobox}

\section{Service Responsibilities}\label{sec:historical-current-background}
The backend is divided into four services, each with its own repository:
\texttt{auth-service}, \texttt{products-service}, \texttt{orders-service}, and
\texttt{payments-service}. Table~\ref{tab:domain-ownership} sets out what each
service owns.

\begin{table}[H]
\centering
\caption{Service responsibilities.}
\label{tab:domain-ownership}
\small
\begin{tabularx}{\textwidth}{>{\raggedright\arraybackslash}p{2.6cm}
                >{\raggedright\arraybackslash}X}
\toprule
\textbf{Service} & \textbf{Owns} \\
\midrule
\texttt{auth-\allowbreak service} & User accounts and admin accounts (separate Mongoose models), credential hashing, admin login tokens \\
\texttt{products-\allowbreak service} & Product catalog and categories; also owns the one-time database seed (pulled from a public demo API, \texttt{api.escuelajs.co}) \\
\texttt{orders-\allowbreak service} & Orders and order line items; validates stock and price and decrements stock at order-creation time \\
\texttt{payments-\allowbreak service} & Stripe Checkout session creation; also creates its own \texttt{Order} record for the checkout flow \\
\bottomrule
\end{tabularx}
\end{table}

\begin{infobox}[Service Communication]
The backend services do not call each other over HTTP. \texttt{orders-service} and
\texttt{payments-service} each declare their own local copy of the
\texttt{product.model.ts} Mongoose schema and query it directly, since all five
workloads connect to the same MongoDB instance and database; Mongoose resolves a
model named \texttt{Product} to the same underlying \texttt{products} collection
regardless of which service defines the schema. The only HTTP orchestration in the
system happens in \texttt{ecomm-ui}'s BFF routes, which call each backend service by
its Kubernetes Service DNS name.
\end{infobox}

\section{Technology Stack}\label{sec:tech-stack}
Table~\ref{tab:tech-stack} summarises the technologies and versions used across the
application and infrastructure layers.

\begin{table}[H]
\centering
\caption{Technology stack.}
\label{tab:tech-stack}
\small
\begin{tabularx}{\textwidth}{>{\raggedright\arraybackslash}p{3.2cm}
                >{\raggedright\arraybackslash}X}
\toprule
\textbf{Layer} & \textbf{Notes} \\
\midrule
Backend runtime & Node.js 20 (\texttt{node:20-slim}), TypeScript run directly via \texttt{tsx} (no separate compile step in \texttt{start}) \\
Backend framework & Express 4.19, \texttt{cors}, \texttt{dotenv} -- identical across all four backend services \\
Backend data access & Mongoose 8.9 against a single shared MongoDB 7.0 instance \\
Auth-specific & \texttt{bcryptjs} 2.4 for password hashing; \texttt{jose} 6.2 (\texttt{SignJWT}, HS256) for admin JWTs only \\
Payments-specific & \texttt{stripe} 17.5 (Node SDK), Stripe Checkout Sessions \\
Frontend & Next.js 15.5 (App Router), React 19, Tailwind CSS 3.4, \texttt{lucide-react} icons \\
Frontend state & React Context + \texttt{localStorage} for the cart (\texttt{CartContext.tsx}); a second Zustand store also exists (\texttt{hooks/useCart.ts}) -- see Section~\ref{sec:frontend-completeness} \\
Frontend dependencies declared but unused & \texttt{zod}, \texttt{next-auth}, \texttt{bcryptjs} (declared in \texttt{package.json}, not currently imported by application code) \\
Container orchestration & Kubernetes v1.34.2 (OKE-managed), Helm 3 chart \texttt{ecommerce-app} \\
Infrastructure as code & Terraform, provider \texttt{oracle/oci} $\geq$ 5.30, provider \texttt{integrations/github} $\sim$ 6.0; remote state in an S3-compatible OCI Object Storage backend \\
GitOps & ArgoCD (upstream manifests, installed by \texttt{deploy.yml}), NGINX Ingress Controller \\
\bottomrule
\end{tabularx}
\end{table}

All five workloads share a single MongoDB instance with no per-service database
split.

\section{Deployed Architecture}\label{sec:deployed-architecture}
Figure~\ref{fig:topology} shows the deployed network and OKE architecture: four
subnets, the Internet/NAT/Service gateway split, and the shared MongoDB PVC in the
Pods subnet.

\begin{figure}[H]
\centering
\resizebox{\textwidth}{!}{%
\begin{tikzpicture}[font=\sffamily\footnotesize]

\node[neutralbox, minimum width=4.8cm, minimum height=1.7cm, text width=4.1cm] (internetN) at (0,16.6)
  {\textbf{Internet}\\[2pt]{\scriptsize shoppers / \texttt{kubectl}}};
\node[neutralbox, minimum width=4.8cm, minimum height=1.7cm, text width=4.1cm] (devN) at (0,13.4)
  {\textbf{Developer}\\[2pt]{\scriptsize\texttt{kubectl} / browser}};
\draw[arrow] (devN.north) -- (internetN.south);

\draw[ThesisGrey, dashed, line width=0.6pt, rounded corners=3pt] (-18.2,-3.6) rectangle (18.2,12.3);
\node[title, anchor=north west] at (-17.8,11.9) {VCN: \texttt{groupa-dev-vcn} (\texttt{10.0.0.0/16})};

\draw[ThesisBlue, line width=0.55pt, rounded corners=1pt, fill=white] (-17.6,5.7) rectangle (-1.0,10.1);
\draw[ThesisBlue, line width=0.55pt, rounded corners=1pt, fill=white] (1.0,5.7) rectangle (17.6,10.1);
\draw[ThesisBlue, line width=0.55pt, rounded corners=1pt, fill=white] (-17.6,-2.9) rectangle (-1.0,3.9);
\draw[ThesisBlue, line width=0.55pt, rounded corners=1pt, fill=white] (1.0,-2.9) rectangle (17.6,3.9);

\node[title, anchor=north west, text width=8.5cm] at (-17.2,9.75) {API Endpoint Subnet (Public)\\[1pt]{\scriptsize\ttfamily 10.0.0.0/28}};
\node[title, anchor=north west, text width=8.5cm] at (1.4,9.75) {Load Balancer Subnet (Public)\\[1pt]{\scriptsize\ttfamily 10.0.0.16/28}};
\node[title, anchor=north west, text width=8.5cm] at (-17.2,3.55) {Worker Nodes Subnet (Private)\\[1pt]{\scriptsize\ttfamily 10.0.16.0/20}};
\node[title, anchor=north west, text width=8.5cm] at (1.4,3.55) {Pods Subnet (Private)\\[1pt]{\scriptsize\ttfamily 10.0.32.0/19}};

\node[evalbox, minimum width=4.0cm, minimum height=1.6cm, text width=3.5cm] (igwN) at (-13.7,7.6) {Internet Gateway};
\node[evalbox, minimum width=4.2cm, minimum height=1.6cm, text width=3.7cm] (okeN) at (-4.5,7.6) {OKE API Endpoint};
\draw[arrow] (igwN.east) -- (okeN.west);

\node[evalbox, minimum width=4.6cm, minimum height=1.6cm, text width=4.1cm] (lbN) at (9.3,7.6) {Load Balancer};

\node[evalbox, minimum width=3.7cm, minimum height=1.5cm, text width=3.2cm] (wn1) at (-13.7,0.5) {Worker Node};
\node[evalbox, minimum width=3.7cm, minimum height=1.5cm, text width=3.2cm] (wn2) at (-9.3,0.5) {Worker Node};
\node[evalbox, minimum width=3.7cm, minimum height=1.5cm, text width=3.2cm] (natN) at (-4.9,0.5) {NAT Gateway};
\draw[flowline] (wn1.east) -- (wn2.west);
\draw[arrow] (wn2.east) -- (natN.west);

\node[evalbox, minimum width=4.0cm, minimum height=1.7cm, text width=3.5cm] (podsN) at (4.9,0.5) {\textbf{Pods}\\[2pt]{\scriptsize 5 services}};
\node[evalbox, minimum width=4.2cm, minimum height=1.7cm, text width=3.7cm] (svcgwN) at (9.3,0.5) {Service Gateway};
\node[outputbox, minimum width=4.0cm, minimum height=1.7cm, text width=3.5cm] (mongoN) at (13.7,0.5) {\textbf{MongoDB}\\[2pt]{\scriptsize PVC}};
\draw[arrow] (podsN.east) -- (svcgwN.west);
\draw[arrow] (podsN.north) -- (4.9,2.5) -- (13.7,2.5) -- (mongoN.north);
\node[note] at (9.3,2.8) {\scriptsize mount};

\node[neutralbox, minimum width=4.4cm, minimum height=1.7cm, text width=3.9cm] (ocirN) at (13.7,-6.9) {OCIR Registry};

\draw[arrow] (internetN.south) -| (igwN.north);
\node[note, text width=3.2cm] at (-10,13.0) {\scriptsize HTTPS/6443};

\draw[arrow] (internetN.south) -| (lbN.north);
\node[note, text width=3.2cm] at (6.5,13.0) {\scriptsize HTTPS/443};

\draw[arrow, dashed] (okeN.south) -- (-4.5,3.9);
\node[note, text width=3.4cm] at (-2.1,4.85) {\scriptsize TCP 6443 /\\12250};

\draw[arrow] (lbN.south) -- (9.3,2.9) -- (-9.3,2.9);
\node[note] at (0,3.2) {\scriptsize routes to};

\node[note] at (-7.1,1.55) {\scriptsize outbound};

\draw[bluearrow] (-1.0,0.5) -- (1.0,0.5);
\node[note] at (0,1.2) {\scriptsize hosts pods};

\draw[arrow] (natN.south) -- (-4.9,-2.3) -- (9.3,-2.3) -- (svcgwN.south);
\node[note] at (2.2,-2.65) {\scriptsize OCI services};

\draw[bluearrow, dashed] (svcgwN.south) -- (9.3,-5.0) -- (ocirN.north);
\node[note] at (11.5,-5.3) {\scriptsize pull image};

\end{tikzpicture}
}
\caption{Deployed architecture inside VCN \texttt{groupa-dev-vcn} (\texttt{10.0.0.0/16}).}
\label{fig:topology}
\end{figure}

Neither the API endpoint path nor the Load Balancer path currently restricts inbound
traffic by source IP; Chapter~\ref{chap:security} covers the network configuration
in full.

Figure~\ref{fig:app-access} shows the application-layer routing: the ingress routes
by path to each of the five workloads, all of which connect to the same MongoDB
instance.

\begin{figure}[H]
\centering
\resizebox{0.85\textwidth}{!}{%
\begin{tikzpicture}[font=\sffamily\footnotesize]
\node[outputbox, minimum width=4.2cm, minimum height=1.6cm, text width=3.8cm] (ingress) at (0,0)
  {\textbf{Ingress}\\{\scriptsize path-based routing}};

\node[evalbox, minimum width=3.1cm, minimum height=1.5cm, text width=2.8cm] (ui) at (-9,-3.4)
  {\textbf{ecomm-ui}\\{\scriptsize\ttfamily port 3000}};
\node[evalbox, minimum width=3.1cm, minimum height=1.5cm, text width=2.8cm] (auth) at (-4.5,-3.4)
  {\textbf{auth-service}\\{\scriptsize\ttfamily port 5001}};
\node[evalbox, minimum width=3.1cm, minimum height=1.5cm, text width=2.8cm] (prod) at (0,-3.4)
  {\textbf{products-service}\\{\scriptsize\ttfamily port 5002}};
\node[evalbox, minimum width=3.1cm, minimum height=1.5cm, text width=2.8cm] (ord) at (4.5,-3.4)
  {\textbf{orders-service}\\{\scriptsize\ttfamily port 5003}};
\node[evalbox, minimum width=3.1cm, minimum height=1.5cm, text width=2.8cm] (pay) at (9,-3.4)
  {\textbf{payments-service}\\{\scriptsize\ttfamily port 5004}};

\node[outputbox, minimum width=4.2cm, minimum height=1.6cm, text width=3.8cm] (mongo) at (0,-6.8)
  {\textbf{MongoDB}\\{\scriptsize shared database}};

% Single distribution bus from Ingress down to each service -- avoids five
% overlapping bends fanning out from one point.
\draw[flowline] (ingress.south) -- (0,-1.7);
\draw[flowline] (-9,-1.7) -- (9,-1.7);
\draw[arrow] (-9,-1.7) -- (ui.north);
\draw[arrow] (-4.5,-1.7) -- (auth.north);
\draw[arrow] (0,-1.7) -- (prod.north);
\draw[arrow] (4.5,-1.7) -- (ord.north);
\draw[arrow] (9,-1.7) -- (pay.north);

% Single collector bus from each service down into MongoDB.
\draw[flowline] (-9,-5.1) -- (9,-5.1);
\draw[flowline] (-9,-4.15) -- (-9,-5.1);
\draw[flowline] (-4.5,-4.15) -- (-4.5,-5.1);
\draw[flowline] (0,-4.15) -- (0,-5.1);
\draw[flowline] (4.5,-4.15) -- (4.5,-5.1);
\draw[flowline] (9,-4.15) -- (9,-5.1);
\draw[arrow] (0,-5.1) -- (mongo.north);
\end{tikzpicture}
}
\caption{Application-layer access structure.}
\label{fig:app-access}
\end{figure}

\begin{infobox}[Note]
\texttt{ecomm-ui} is not a pure stateless proxy in front of the backend services.
Most of its \texttt{/api/*} routes forward requests to the corresponding backend
service. The frontend's actual data path is HTTP calls to the four backend
services.
\end{infobox}

\texttt{Terraform-code} includes two deployment definitions for the same
application: the Helm chart (\texttt{helm/ecommerce-app}) and a parallel set of raw
manifests (\texttt{kubernetes/00} through \texttt{06}). ArgoCD's
\texttt{Application} resource points at \texttt{helm/ecommerce-app}, so the Helm
chart is the path GitOps manages; the raw manifests are a second, manually applied
option documented in \texttt{DEPLOYMENT.md} as ``Option B'' and are not kept in sync
with the Helm chart automatically. The two paths currently define different
autoscaling policies (Section~\ref{sec:scaling-resource-allocation}).

\section{Network Architecture}\label{sec:domain-context}
The OKE cluster network uses four subnets, summarised in
Table~\ref{tab:subnet-design}. The full set of gateway, route table, and NSG rules
is in Appendix~\ref{sec:appendix-network}.

\begin{table}[H]
\centering
\caption{OKE subnet configuration.}
\label{tab:subnet-design}
\small
\begin{tabularx}{\textwidth}{>{\raggedright\arraybackslash}p{2.4cm} l
                >{\raggedright\arraybackslash}X l l}
\toprule
\textbf{Subnet} & \textbf{Type} & \textbf{Purpose} & \textbf{CIDR} & \textbf{Range} \\
\midrule
API Endpoint & Public & Kubernetes API access & 10.0.0.0/28 & .0--.15 \\
Load Balancer & Public & Public entry point for the application & 10.0.0.16/28 & .16--.31 \\
Worker Nodes & Private & Hosts the VMs running all pods & 10.0.16.0/20 & 10.0.16.0--10.0.31.255 \\
Pods & Private & IP addresses for individual pods (VCN-native) & 10.0.32.0/19 & 10.0.32.0--10.0.63.255 \\
\bottomrule
\end{tabularx}
\end{table}

\begin{infobox}[Note]
The network module (\texttt{modules/network/variables.tf}) defines four boolean
switches, each defaulting to \texttt{false}, that control worker SSH access,
general worker outbound traffic, pod outbound internet access, and public
NodePort access. The root module (\texttt{main.tf}) sets all four to
\texttt{true}. Chapter~\ref{chap:security} lists exactly what this opens and
covers the security implications.
\end{infobox}

\section{Scaling and Routing}\label{sec:app-layer-design}
Every workload scales through its own Horizontal Pod Autoscaler, and the ingress
routes to each service by path. Section~\ref{sec:scaling-resource-allocation} covers
the exact policy and a difference between the two deployment paths.

\section{Design Draft vs.\ Implementation}\label{sec:draft-vs-implementation}
The original Architecture \& Design Draft set out the intended design before
implementation began. Table~\ref{tab:draft-vs-impl} summarises how the current
implementation compares, including areas the draft did not originally cover.

\begin{table}[H]
\centering
\caption{Design draft vs.\ current implementation.}
\label{tab:draft-vs-impl}
\small
\begin{tabularx}{\textwidth}{>{\raggedright\arraybackslash}p{3.6cm}
                >{\raggedright\arraybackslash}p{2.6cm}
                >{\raggedright\arraybackslash}X}
\toprule
\textbf{Area} & \textbf{Classification} & \textbf{Current implementation} \\
\midrule
Domain-to-service mapping & Implemented & Four domains map 1:1 to four repositories, as designed \\
Inter-service communication & Implemented differently & The original design anticipated service-to-service API calls; the implemented services instead access shared MongoDB collections directly (Section~\ref{sec:historical-current-background}) \\
NSG-restricted API/SSH access & Implemented differently & The root Terraform configuration enables all four optional network access switches and allows public access to the Kubernetes API endpoint (Section~\ref{sec:domain-context}) \\
Secrets via OCI Vault & Planned & Referenced in \texttt{DEPLOYMENT.md} but not implemented in the current Terraform configuration or GitHub Actions workflows (Chapter~\ref{chap:security}) \\
Role-based product/order authorization & Planned & Backend services and BFF routes do not currently check a JWT or role before mutating data (Chapter~\ref{chap:security}) \\
User authentication & Partially implemented & Credentials are validated server-side, but no session or token is currently issued to a regular user (Chapter~\ref{chap:security}) \\
Payment confirmation & Partially implemented & A Stripe Checkout session is created server-side; order status is confirmed through a client call rather than a webhook or session check (Chapter~\ref{chap:security}) \\
Full storefront browsing & Partially implemented & The homepage and admin panels are functional; the product listing, product detail, and category pages remain incomplete (Section~\ref{sec:frontend-completeness}) \\
Two parallel deployment manifests & Repository-only & A raw-manifest path and a Helm-chart path both exist and currently define different HPA policies; the draft describes neither (Section~\ref{sec:scaling-resource-allocation}) \\
Automated testing & Repository-only & Automated tests are not included in any of the six repositories \\
\bottomrule
\end{tabularx}
\end{table}
